\section{Datasets and Models}
\label{sec:dataset}

\paragraph{Datasets.} We evaluate six datasets, each filtered to items where judges are more likely to be uncertain (Appendix~\ref{sec:appendix-sampling}).
The safety tasks include \textbf{WildGuard} \citep{wildguardmix2024} (adversarial prompts with compliant responses),
\textbf{AEGIS} \citep{aegis2024} (a second safety taxonomy for replication), and \textbf{HH-RLHF} \citep{ganguli2022redteam} 
(Anthropic's red-team-attempts, stratified across harm levels 0--4). The remaining datasets are \textbf{ToxiGen} 
\citep{hartvigsen2022toxigen} (adversarial toxicity items balanced across demographic targets), \textbf{MAGE} 
\citep{li2024mage} (AI-generated vs.\ human-written text detection with known provenance), and \textbf{Paired Prompts} 
\citep{anthropicPoliticalNeutrality} (political content with \emph{hedging} and \emph{refusal}rubrics).
Sample sizes are 100 items per safety/toxicity/AI-detection dataset and 50 prompt pairs per Paired Prompts rubric. 
Five of the datasets have ground-truth labels, which allows us to classify each wiggle as \emph{corrective} or \emph{corrupting}.
Corrective wiggles move the judge's L0 verdict toward the ground-truth label, while corrupting wiggles move it away.
Paired Prompts has no canonical ground truth and is excluded from corrective/corrupting analysis.
Full per-task wiggle rates, a compact task summary, and model-by-task retention rates are reported in
Tables~\ref{tab:wiggle-rates}, \ref{tab:task-summary}, and~\ref{tab:model-retention}, respectively.

\paragraph{Judging tasks.} Each judging task is tested with a Binary scale and a 1-5 Likert scale, giving 14 (dataset, rubric, scale) judging tasks in total (5 single-rubric datasets $\times$ 2 scales + 1 dataset (Paired Prompts) $\times$ 2 rubrics $\times$ 2 scales).
Verdict spaces are organized into \emph{restrictive} and \emph{permissive} sides. A flip is considered restrictive if the judge 
takes a more conservative action. On the safety datasets (WildGuard, AEGIS, HH-RLHF), 
restrictive is \emph{unsafe / harmful} and permissive is \emph{safe / helpful}. On ToxiGen, restrictive is \emph{toxic} 
and permissive is \emph{benign}. On MAGE, restrictive is \emph{AI-generated} (the suspicion side, parallel to 
\emph{unsafe} on safety) and permissive is \emph{human-written}. Two different rubrics are used for Paired Prompts: 
on the hedging rubric, restrictive is \emph{more hedging}\footnote{When judging non-anchored political prompts, it's plausible to
argue that \emph{direct response} is the more restrictive position. Since we use the same judge system prompt as the original benchmark, 
which goes from most-direct (1) to most-hedging (5), we maintain parallelism with other tasks where 1 $\rightarrow$ 5 goes from most permissive $\rightarrow$ most restrictive.}. 
On the refusal rubric, the restrictive side is \emph{more refusing} and permissive is \emph{compliant}.

\paragraph{Models.} We evaluate 9 judge models across four families: GPT-5, GPT-5.2, GPT-5.4 (OpenAI); 
Claude 4.6 Sonnet, Claude 4.6 Opus (Anthropic); Grok-4.1, Grok-4.1 Reasoning (xAI); Gemini 3 Flash, Gemini 3.1 Pro (Google). 
For L6 adaptive persuasion, three of these (GPT-5.4, Claude Opus, Grok-4.1 Reasoning) serve as persuaders, generating 
challenges for all judges including themselves. To check whether the judge's verdict has flipped, we use an observer model 
(GPT-5) to parse the judge's verdict from free-form responses. Where a temperature parameter is accepted, models are queried
at temperature=0, including for the L0 baseline, and default temperature otherwise. We use each model's default reasoning 
configuration.\footnote{Reasoning is off for OpenAI's models and for the 
Claude models, extended thinking is disabled. Reasoning is on for Grok-4.1 Reasoning and for both Gemini 3 models with 
default reasoning effort.}

% \paragraph{Jury baseline.} We additionally compute the jury's L0 majority strength, following \citet{zhao2024council}. 
% The jury's majority strength (the fraction of judges agreeing with the modal L0 verdict) for the majority verdict is 
% included as a feature in the cross-level correlation analysis (\S\ref{sec:discussion}).
